\begingroup
\makeatletter\tagsleft@true\let\veqno\@@leqno\makeatother
\ArithmeticHead{280}{L}{ANALYSIS. NACHLASS.}
\SourceLeaf{287b}
\pdfbookmark[1]{Erläuterungen zu [IX.]}{x1-agm-erlaeuterung-ix}
\begin{center}\textls[120]{Erläuterungen zu [IX.], S. 217—231.}\end{center}
\par\indent In \textsc{Scherings} Untersuchungen über das agM., die Werke III, S. 374 abgedruckt sind, finden sich zwei Gleichungen, die eine S. 377 Zeile 3, die andere S. 380 Zeile 13, 14, von denen \textsc{Schering} S. 380 sagt, daß \textsc{Gauss} sie nebeneinander aufgezeichnet habe. In der Tat sind diese Gleichungen im art. [1.] des
\clearpage
\fancyhead{}\fancyhead[R]{\small281}\fancyhead[C]{\scriptsize BEMERKUNGEN ZUM AGM. ABSCHNITT IX.}
\SourceLeaf{288}
\noindent Abschnitts [IX.] enthalten, und zwar stimmt die erste mit der Gl. [1], die zweite mit der Gl. [4] überein. Die Bedeutung namentlich der ersten Gleichung erhellt aus ihrem Zusammenhang mit dem »Schönen Lehrsatz« des art. [4.], S. 220, vergl. den Artikel 2. unseres \textit{Abrisses} oben S. 255 und die Erörterungen von \textsc{Schering} im Artikel 13. seines Aufsatzes, Werke III, S. 379, 380; die Gleichung der ersten Zeile von S. 380 ist nichts anderes als jener »Schöne Lehrsatz«, der also auch mit dem Werke VIII, S. 98 abgedruckten \textit{Theorema elegantissimum} übereinstimmt (vergl. die Bemerkung von R. \textsc{Fricke}, a. a. O.). Übrigens ist die Gleichung [1] vollkommen identisch mit der historisch so bedeutungsvollen Gleichung [7] des Abschnitts [IV.] (oben S. 186), die in der Scheda Ac, S. 10 aufgezeichnet ist. In [1], [2], [3] bedeutet $k$ den Modul der \textsc{Brigg}schen Logarithmen, vergl. auch art. [2.], S. 218, wo die mit $a$ bezeichnete Größe gleich $k\pi$ ist. Die Handschrift enthält noch eine Formel für die Peripheria Ellipsis, die mit der in der Anzeige der \textit{Determinatio attractionis} (Werke III, S. 360) gegebenen übereinstimmt und darum hier nicht abgedruckt worden ist.
\par\indent Die im art. [1.] enthaltene Aufzeichnung stammt aus dem Mai 1809.
\par\indent Die in den artt. [2.]—[9.] wiedergegebenen Aufzeichnungen geben eine vollständige Skizze für die Theorie der Modulfunktion. Der Inhalt der artt. [2.] und [4.] ist im wesentlichen in der Scheda Af (siehe die artt. [2.], [3.], [4.] des Abschnitts [VII.], S. 210, 211) enthalten. Der Satz [1] des art. [3.], S. 219 gibt den \textit{nexum mutuum}, der die \textit{infinite multos terminos medii arithmetico-geometrici} mit einander verbindet, und den \textsc{Gauss} nach der Notiz Nr. 109 des \textit{Tagebuchs} am 3. Juni 1800 in Braunschweig gefunden hat. Zu diesem Fundamentalsatz der Lehre von der Modulfunktion ist folgendes zu bemerken.
\par\indent Bedeutet $\mathfrak M(a,b)$ einen \textls[120]{beliebigen} von Null verschiedenen Wert des agM. zwischen $a$ und $b$, bei dessen Erzeugung durch den Algorithmus des agM. die Vorzeichen der Quadratwurzeln also irgendwie gewählt worden sind, so ist stets
\begin{equation*}\tag*{$(\ast)$}\frac1{\mathfrak M(a,b)}=\frac\delta{M(a,b)}+\frac{i\gamma}{M(a,c)}\mathpunct{\raisebox{.5ex}{.}}\end{equation*}
\noindent Wenn in dieser Formel $M(a,b),\ M(a,c)$ \textls[120]{irgend} zwei bestimmte Werte des agM. zwischen $a$ und $b$, beziehungsweise zwischen $a$ und $c=\sqrt{a^2-b^2}$ bedeuten, so sind $\delta,\gamma$ ganze Zahlen, die den Bedingungen
\[\delta\equiv1,\quad\gamma\equiv0\ (\text{mod. }4)\]
\noindent genügen. Nennt man zwei Werte $M(a,b),\ M(a,c)$ \textls[120]{zusammengehörig}, wenn der reale Teil des Quotienten $\frac{M(a,b)}{M(a,c)}$ wesentlich positiv ist, so gilt\textsuperscript{*)} das folgende. Bedeuten $M(a,b),\ M(a,c)$ zwei zusammengehörige Werte, so ist in der Formel $(\ast)$ dann und nur dann $\delta=1$, wenn auch $\mathfrak M(a,b)$ und $M(a,c)$ zusammengehörige Werte sind. Das ist also die Bedeutung der \textsc{Gauss}schen Formel [1] oder\textsuperscript{**)}
\begin{equation*}\tag*{$[1]'$}\frac1{\mathfrak M(a,b)}=\frac1{M(a,b)}+\frac{4ik}{M(a,c)}\mathpunct{\raisebox{.5ex}{.}}\end{equation*}
\par\indent Wahrscheinlich hat \textsc{Gauss} diese Formel aus der Theorie der Reihen $p,q,r$, die er als die \textls[120]{Summatorischen Funktionen} bezeichnet\textsuperscript{***)}, abgeleitet; diese Vermutung stützt sich auf den art. [6.], S. 222, wo in [5] der allgemeinste Wert der dort mit $M$ bezeichneten Größe angegeben und in den Kettenbruch [4] entwickelt wird, bei dem der Quotient
\[\frac nm=\frac{q^2}{p^2}\]
\Needspace{4\baselineskip}
\par\smallskip\noindent\rule{.25\linewidth}{.4pt}\par
{\small
\par\indent *) Nach einer Mitteilung von L. v. \textsc{Dávid}.
\par\indent **) Vergl. Werke III, S. 378.
\par\indent ***) Werke III, S. 386 unten; diese Stelle, d. h. die fünf letzten Zeilen dieser Seite, hat \textsc{Gauss} auf der letzten Einbandseite seines Handexemplars der \textit{Disquisitiones arithmeticae} aufgezeichnet.
\par}
\par\smallskip\noindent X1.\hfill36\par
\clearpage
\ArithmeticHead{282}{L}{ANALYSIS. NACHLASS.}
\SourceLeaf{289}
\noindent seinen Wert nicht verändert. Hier ist also
\[M=\frac{M(m,n)}{M(m,\sqrt{m^2-n^2})};\]
\noindent auch wendet \textsc{Gauss} in den artt. [6.], [7.], [8.], [9.] die Funktionszeichen $p,q,r$ etwas abweichend von der sonst gebrauchten Weise an, indem er nämlich
\[\begin{aligned}
p(t)&=1+2e^{-\pi t}+2e^{-4\pi t}+2e^{-9\pi t}+\ldots,\\
q(t)&=1-2e^{-\pi t}+2e^{-4\pi t}-2e^{-9\pi t}+\ldots,\\
r(t)&=2e^{-\frac14\pi t}+2e^{-\frac94\pi t}+2e^{-\frac{25}4\pi t}+\ldots
\end{aligned}\]
\noindent setzt; vergl. die Werke VIII, S. 101 in art. [3.] abgedruckte Aufzeichnung, die aus dem Handbuch 25, Bh, S. 25 stammt und etwa 1825 geschrieben sein dürfte. Daß auf S. 223 mit den Formeln für die lineare Transformation von $pt$ auch die entsprechenden Formeln für die von den \textls[120]{beiden} Veränderlichen abhängige Thetafunktion aufgezeichnet sind (vergl. die ausführliche Darstellung dieser Theorie im Abschnitt [I.] artt. [1.] und [13.] der unten folgenden Abhandlung \textit{Zur Theorie der transscendenten Functionen gehörig}, besonders die letzte Gleichung des art. [13.]), erscheint durch die Form bemerkenswert, in der \textsc{Gauss} hier die allgemeine Thetafunktion schreibt, nämlich
\[\sum_n h^{\alpha n^2+2\beta n+\gamma},\]
\noindent wo die quadratische binäre Form im Exponenten den Zusammenhang mit der Zahlentheorie unmittelbar hervortreten läßt.
\par\indent In den artt. [7.] und [8.] wird das Verhalten der drei Summatorischen Funktionen $p,q,r$ bei ganzzahliger linearer Transformation von $M$ in den sechs verschiedenen Fällen modulo 2 aufgestellt und mit Zuhilfenahme der Reduktionstheorie der binären quadratischen Formen negativer Determinante, die »einfachste Form« von $M$ gefunden. Daß in der Tabelle [6], S. 224, für $h$ der reziproke Wert des in der Handschrift und Werke III, S. 386 angegebenen Wertes $i^\alpha\surd\delta+\gamma t)$ genommen werden muß, hat schon P. \textsc{Pepin}\textsuperscript{*)} bemerkt. Vergl. im übrigen den Abschnitt V. des Aufsatzes »Über \textsc{Gauss}' Arbeiten zur Funktionentheorie«.
\par\indent Die beiden Zettel, auf denen die art. [1.]—[9.] aufgezeichnet sind, stammen aus verhältnismäßig später Zeit; das geht u. a. daraus hervor, daß auf dem zweiten Zettel (siehe oben S. 226) die von \textsc{Fourier} 1822 eingeführte Bezeichnung der Grenzen eines bestimmten Integrals benutzt wird. Sie dürften etwa 1825 geschrieben sein, aber, wie bereits bemerkt, stammt ihr wesentlicher Inhalt aus dem Jahre 1801. Auszüge aus diesen beiden Zetteln sind Werke III, S. 378, 379, 385—386 und VIII, S. 99, 100—101 abgedruckt.
\par\indent In den artt. [10.] und [12.] wird die Darstellung des elliptischen Integrals erster und zweiter Gattung mit veränderlicher oberer Grenze ähnlich wie in der \textit{Determinatio attractionis} gegeben; die Rektifikation der Ellipse wird im art. [13.] noch nach einer zweiten Methode behandelt. Die Entstehungszeit dieser Untersuchungen ist durch die Tagebuchnotizen Nr. 110 und Nr. 111 auf 5. bis 10. Juni 1800 festgelegt. Die \textit{Transcendentes ellipticae}, auf die nach Nr. 110 \textsc{Gauss} seine Theorie nun auch \textls[120]{unmittelbar} angewendet hat, sind die Integrale erster Gattung, und mit der unmittelbaren Anwendung ist der \textit{Algorithmus} gemeint, der $\theta=T^\infty$ liefert (siehe oben S. 228). Von den drei verschiedenen Methoden, nach denen \textsc{Gauss} gemäß der Notiz Nr. 111 die Rektifikation erledigt hat, sind hier zwei entwickelt. Die in den artt. [10.]—[14.] abgedruckten drei Zettel selbst sind ersichtlich kurz hintereinander, aber auch sehr viel später als 1800
\Needspace{4\baselineskip}
\par\smallskip\noindent\rule{.25\linewidth}{.4pt}\par
{\small
\par\indent *) \textit{Introduction à la théorie des fonctions elliptiques d'après les Oeuvres posthumes de Gauss}, Memorie dell’ Accad. Pontif. de’ Nuovi Lincei \textbf{9}, 2, 1893, S. 103.
\par}
\clearpage
\fancyhead{}\fancyhead[R]{\small283}\fancyhead[C]{\scriptsize BEMERKUNGEN ZUM AGM. BRIEFWECHSEL.}
\SourceLeaf{290a}
\noindent geschrieben\textsuperscript{*)}. Der im art. [13.], S. 230 enthaltenen Rechnung liegt nämlich der \textsc{Walbeck}sche Wert der Abplattung der Erde $\frac1{302{,}78}$ zu Grunde, der aus dem Jahre 1819 stammt\textsuperscript{**)} und mit dem \textsc{Gauss} in den zwanziger Jahren des XIX. Jahrhunderts gearbeitet hat\textsuperscript{***)}. Bezeichnet man mit $\mathrm a$ die halbe große, mit $\mathrm b$ die halbe kleine Achse der Meridianellipse, so ist
\[a=302{,}78=\frac{\mathrm a}{\mathrm a-\mathrm b},\quad b=1-a=\frac{\mathrm b}{\mathrm a-\mathrm b},\]
\noindent also in der Tat
\[\frac1a=\frac{\mathrm a-\mathrm b}{\mathrm a}=\frac1{302{,}78}\]
\noindent die Abplattung. Der \textit{factor constans} ist also $(\mathrm a-\mathrm b)^{-1}$, und der sich ergebende Wert $\log_{10}\mathrm a=6{,}8045975$ stimmt mit dem Werke IX, S. 69 berechneten (dort ist die halbe große Achse mit $a$ bezeichnet) überein; zur Berechnung des Ellipsenumfangs dient dort die gewöhnliche Reihe, siehe im \textsc{Leiste}, oben S. 177, Gl. [5].
\Needspace{4\baselineskip}
\par\smallskip\noindent\rule{.25\linewidth}{.4pt}\par
{\small
\par\indent *) Der erste, der den art. [10.] enthält, ist anscheinend von einem Buchdeckel abgelöst; die beiden andern und noch ein dritter, mit Zahlenrechnungen und der Wiederholung von Formeln, die auch auf den andern stehen, sind von demselben Bogen abgetrennt. In der Formel des art. [10.], S. 227 Zeile 9, enthält die Handschrift ein Versehen, das Integral im dritten Gliede der Gleichung muß nämlich lauten
\[\int\frac{dT''}{\surd(a''a''\sin4T''^2+b''b''\cos4T''^2)};\]
\noindent dabei ist natürlich
\[a'=\frac{a+b}2,\quad b'=\sqrt{ab},\quad a''=\frac{a'+b'}2,\quad a''=\sqrt{a'b'}.\]
\par\indent **) \textsc{Henric. Joh. Walbeck}, \textit{Dissertatio de forma et magnitudine telluris etc.} Aboae 1819.
\par\indent ***) Vergl. die \textit{Zahlen, das Erdsphäroid betreffend}, abgedruckt Werke IX, S. 69 und die zugehörige Bemerkung ebenda S. 71, ferner die Briefe von \textsc{Gauss} an \textsc{Olbers} vom 18. April 1822, 6. Juli 1824, 1. März 1827, abgedruckt ebenda S. 369, 320, 377.
\par}
\clearpage
\endgroup
